\Article{5}
Formula, per quam in art.\ praec.\ productum e functionibus quotcunque expressimus, generalis est, neque ullum nexum inter radices $r,r',r'',r'''$ etc., vel inter $R,R',R'',R'''$ etc.\ supponit. Nullo inde negotio deducitur, si radices $r',r'',r'''$ etc.\ tamquam potestates ipsius $r$, radicesque $R',R'',R'''$ etc.\ tamquam potestates ipsius $R$ considerare liceat, singulas partes producti sub forma $R^M[r^m,R^\lambda]$ comprehensas fore, ubi exponens $\lambda$ pro singulis idem erit, scilicet $R^\lambda=RR'R''R'''$ etc. Quamobrem per ea, quae in art.\ 3 monuimus, huiusmodi productum reducetur ad formam sequentem
\[
\begin{aligned}
&A[1,R^\lambda]+B[r,R^\lambda]+B'[r^g,R^\lambda]+B''[r^{g^2},R^\lambda]+B'''[r^{g^3},R^\lambda]+\text{etc.}\\
&\hspace{5em}{}+B^{(\alpha-1)}[r^{g^{\alpha-1}},R^\lambda]
\end{aligned}
\]
ubi singuli coëfficientes $A,B,B',B'',B'''$ etc.\ erunt formae
\[
h+h'R+h''R^2+h'''R^3+\text{etc.}+h^{(\beta-1)}R^{\beta-1}
\]
designantibus $h,h',h'',h'''$ etc.\ numeros determinatos integros.

\SourceLeaf{257}
Casus simplicissimus is erit, ubi ponitur $r=r'=r''=r'''$ etc., nec non $R=R'=R''=R'''$ etc.; tunc productum nostrum transit in potestatem $[r,R]^\lambda$, quae itaque ad formam supra traditam semper reveniet.

\Article{6}
Statuendo itaque $\lambda=\beta$, potestas $[r,R]^\beta$ hanc formam nanciscetur:
\[
\begin{aligned}
&A[1,1]+B[r,1]+B'[r^g,1]+\text{etc.}+B^{\alpha-1)}[r^{g^{\alpha-1}},1]\\
&{}=\beta\gamma A+B(\beta\gamma,1)+B'(\beta\gamma,g)+B''(\beta\gamma,g^2)+\text{etc.}\\
&\qquad{}+B^{(\alpha-1)}(\beta\gamma,g^{\alpha-1})=\theta'
\end{aligned}
\]

Quodsi itaque non modo valor radicis $R$ (adeoque et valores coëfficientium $A,B,B'$ etc.), sed etiam valores singulorum aggregatorum $\beta\gamma$ terminorum $(\beta\gamma,1),(\beta\gamma,g)$ etc.\ cogniti supponuntur, valor ipsius $\theta'$ sponte innotescet, unde erui poterit $[r,R]$ per formulam ${}^{\beta}\!\surd\theta'$. Haec expressio $\beta$ valores diversos admittit; unde dubium videri posset, quemnam adoptare oporteat: facile autem ostenditur, hoc prorsus arbitrarium esse, quoties $R$ sit radix \emph{propria} aequationis $x^\beta-1=0$. In hoc enim casu patet, illos $\beta$ valores expressionis radicalis ${}^{\beta}\!\surd\theta'$ fore
\[
[r,R],\quad[r^{g^\alpha},R],\quad[r^{g^{2\alpha}},R]\dotsfour[r^{g^{\alpha\beta-\alpha}},R]
\]
quippe quarum functionum potestates $\beta^{\mathrm{tae}}$ per art.\ 3 inter se aequales erunt, ipsae vero inter se ipsis $\beta$ radicibus diversis aequationis $x^\beta-1=0$ proportionales: sed quamdiu aggregata $\beta\gamma$ terminorum $(\beta\gamma,1),(\beta\gamma,g)$ etc.\ tantum cognita sunt, ipsa radix $r$ eatenus tantum determinata est, quod in complexu $(\beta\gamma,1)$ contenta esse debet, arbitrariumque manet, quamnam ex hoc complexu pro $r$ adoptemus. Hae radices vero sunt $r,r^{g^\alpha},r^{g^{2\alpha}}$ etc., et proin etiam e functionibus $[r,R],[r^{g^\alpha},R],[r^{g^{2\alpha}},R]$ etc.\ quamlibet pro $[r,R]$ adoptare possumus.

Hae conclusiones non valerent, si $R$ non esset radix propria aequationis $x^\beta-1=0$: supponendo enim, $R$ esse radicem propriam aequationis $x^{\beta'}-1=0$, ita ut $\beta'$ sit divisor ipsius $\beta$, facile patet, fieri
\[
[r,R]=[r^{g^{\alpha\beta'}},R],\quad[r^{g^\alpha},R]=[r^{g^{\alpha\beta'+\alpha}},R]\text{ etc.}
\]
adeoque in complexu $\beta$ functionum $[r,R],[r^{g^\alpha},R]\ldots[r^{g^{\alpha\beta-\alpha}},R]$ tantummodo $\beta'$ diversas reperiri, et proin etiam e valoribus expressionis ${}^{\beta}\!\surd\theta'$ haud plures quam $\beta'$ admissibiles esse, reliquos $\beta-\beta'$ autem spurios. At nullo negotio perspicitur, in hoc casu haud opus esse usque ad potestatam $\beta^{\mathrm{tam}}$ functionis $[r,R]$ 
\SourceLeaf{258}
ascendere, sed iam potestatem $[r,R]^{\beta'}$ ad formam nostram
\[
\beta\gamma A+B(\beta\gamma,1)+B'(\beta\gamma,g)+B''(\beta\gamma,g^2)\text{ etc.}
\]
reduci. Habebimus itaque $[r,R]$ per expressionem talem ${}^{\beta'}\!\surd\theta$, nihilque intererit, quemnam valorem huius expressionis adoptemus.

\Article{7}
Perinde ut $[r,R]$ etiam functiones $[r,R^2]$, $[r,R^3]$ etc.\ sive generaliter $[r,R^k]$ determinare licebit: patet enim, si substituendo in $\theta'$ loco ipsius $R$ potestates $R^2,R^3$ etc.\ $R^k$ emergere supponantur functiones $\theta'',\theta'''$ etc.\ $\theta^{(k)}$, fore $[r,R^2]^\beta=\theta''$, $[r,R^3]^\beta=\theta'''$ etc.\ et generaliter $[r,R^k]^\beta=\theta^{(k)}$; quamobrem hae quoque functiones per expressiones radicales exprimi poterunt, $[r,R^2]={}^{\beta}\!\surd\theta''$ etc. Sed haud convenit, hisce expressionibus radicalibus uti, quoties quantitas aliqua per functionem ipsarum $[r,R]$, $[r,R^2]$ etc.\ exprimenda est. Scilicet quum singularum valores haud penitus determinati sint, dubium maneret, quosnam inter se combinare liceret: manifesto autem hoc neutiquam arbitrarium est; facile enim perspicitur, simulac pro $[r,R]$ valor determinatus accipiatur, etiam omnes $[r,R^2]$, $[r,R^3]$ etc.\ valores penitus determinatos nancisci debere, qui autem per expressiones radicales non indicantur. His itaque reiectis, expressiones alias indagare oportet, quarum adiumento $[r,R^2]$, $[r,R^3]$ etc.\ \emph{rationaliter} per $[r,R]$ atque quantitates cognitas exhibeantur, quod facile sequenti modo efficimus.

Per theorema art.\ 4, eaque quae in art.\ 5 docuimus, etiam productum $[r,R^k]\times[r,R]^{\beta-k}$ ad formam talem
\[
\begin{aligned}
&\beta\gamma A+B(\beta\gamma,1)+B'(\beta\gamma,g)+B''(\beta\gamma,g^2)+\text{etc.}\\
&\qquad{}+B^{(\alpha-1)}(\beta\gamma,g^{\alpha-1})
\end{aligned}
\]
reducetur, ubi $A,B,B',B''$ etc.\ erunt functiones rationales ipsius $R$. Positis itaque productis
\[
\begin{aligned}
[r,R^2]\times[r,R]^{\beta-2}&=\vartheta''\\
\relax[r,R^3]\times[r,R]^{\beta-3}&=\vartheta'''\\
\relax[r,R^4]\times[r,R]^{\beta-4}&=\beta''''\\
&\text{etc.}
\end{aligned}
\]
erunt etiam $\vartheta'',\vartheta''',\vartheta''''$ etc.\ quantitates rationaliter assignabiles, atque
\SourceLeaf{259}
\[
\begin{aligned}
[r,R^2]&=\frac{\vartheta''}{\theta'}[r,R]^2\\
\relax[r,R^3]&=\frac{\vartheta'''}{\theta'}[r,R]^3\\
\relax[r,R^4]&=\frac{\vartheta''''}{\theta'}[r,R]^4\\
&\text{etc.}
\end{aligned}
\]

Hae expressiones itaque valores functionum $[r,R^2]$, $[r,R^3]$ etc.\ rationaliter exhibent, siquidem non fuerit $[r.R]=0$, in quo casu indeterminatae fierent: at rigorose demonstrare possumus, numquam fieri posse $[r,R]=0$, quoties quidem $r$ denotet radicem ab $1$ diversam, etiamsi expositionem huius demonstrationis, ne hic nimis prolixi fiamus, ad aliam occasionem nobis reservare oporteat.

\Article{8}
Quae in artt.\ praec.\ exposuimus, usum praestant, si a periodis $\beta\gamma$ terminorum ad periodos $\gamma$ terminorum descendere propositum est. Nullo scilicet negotio perspicitur, denotante $R$ radicem propriam, haberi
\[
\begin{aligned}
\beta(\gamma,1)&=(\beta\gamma,1)+[r,R]+[r,R^2]+[r,R^3]+\text{etc.}+[r,R^{\beta-1}]\\
\beta(\gamma,g^\alpha)&=(\beta\gamma,1)+R^{\beta-1}[r,R]+R^{\beta-2}[r,R^2]\\
&\qquad{}+R^{\beta-3}[r,R^3]+\text{etc.}+R[r,R^{\beta-1}]\\
\beta(\gamma,g^{2\alpha})&=(\beta\gamma,1)+R^{2\beta-2}[r,R]+R^{2\beta-4}[r,R^2]\\
&\qquad{}+R^{2\beta-6}[r,R^3]+\text{etc.}+R^2[r,R^{\beta-1}]\\
&\text{etc.}
\end{aligned}
\]

Si hic pro singulis $[r,R]$, $[r,R^2]$ etc.\ expressiones radicales ${}^{\beta}\!\surd\theta'$, ${}^{\beta}\!\surd\theta''$ etc.\ acciperentur, valor cuiusvis seriei inter valores $\beta^{\beta-1}$ dubius esset, qui contra adoptatis expressionibus rationalibus pro $[r,R^2]$ etc.\ ambiguitati alii non erit obnoxius, nisi quae per rei naturam est inevitabilis. Haec observatio attentionem ill.\ \textsc{Lagrange} subterfugisse videtur, qui methodum nostram in Disquis.\ arithm.\ art.\ 360 traditam, ubi haud inconsulto neglectis expressionibus radicalibus solas rationales proposueramus, \emph{simplificavisse} sibi visus est, dum illas pro his substituit (Traité de la résolution numérique des équations; édition $2^{\mathrm{me}}$ pag.\ 311).

Ceterum vix opus est hic monere, simulac valores periodorum $(\gamma,1)$, $(\gamma,g^\alpha)$ etc., aut tantummodo unius ex ipsis eruti sint, valores omnium reliquarum periodorum $\gamma$ terminorum rationaliter inde deduci posse. Descensus itaque a periodis $\beta\gamma$ terminorum ad periodos $\gamma$ terminorum requirit solutionem aequationum $x^\beta=1$, $x^\beta=\theta'$, operationesque reliquae rationaliter perficiuntur.
